% Source: printed p. 238; PDF p. 247
\paragraph{Lời giải 5.1.2.39.}
\begin{enumerate}[label=\alph*.]
\item Chỉ có các cấu xạ đồng nhất, mỗi đối tượng có một cấu xạ: $\{\operatorname{id}_1,\operatorname{id}_2,\operatorname{id}_3,\operatorname{id}_4\}$.
\item Một hàm tử $F:\mathcal D_3\to\mathcal D_2$ gồm một hàm $F:\operatorname{Ob}(\mathcal D_3)\to\operatorname{Ob}(\mathcal D_2)$ trên đối tượng, cùng với một hàm trên cấu xạ bảo toàn đồng nhất và hợp thành. Nhưng vì các cấu xạ duy nhất trong $\mathcal D_3$ là các đồng nhất, phần trên cấu xạ không có sự lựa chọn nào và cũng không có ràng buộc nào. Nói cách khác, một hàm tử $\mathcal D_3\to\mathcal D_2$ được xác định hoàn toàn bởi một hàm $\underline{3}\to\underline{2}$. Có tám hàm như vậy; theo (2.5), có thể ký hiệu chúng bằng các dãy:
\[\begin{gathered}(1,1,1);\ (1,1,2);\ (1,2,1);\ (1,2,2);\\(2,1,1);\ (2,1,2);\ (2,2,1);\ (2,2,2).\end{gathered}\]
\end{enumerate}\hfill$\diamond$
\paragraph{Bài tập 5.1.2.40.} Cho $\mathcal C$ là một phạm trù. Có bao nhiêu hàm tử $\mathcal C\to\mathcal D_1$, trong đó $\mathcal D_1:=\operatorname{Disc}(\underline{1})$ là phạm trù rời rạc trên một phần tử?\hfill$\diamond$
\paragraph{Lời giải 5.1.2.40.} Luôn có đúng một hàm tử $\mathcal C\to\mathcal D_1$. Không có lựa chọn nào về nơi gửi các đối tượng (tất cả đều đến đối tượng $1$), và không có lựa chọn nào về nơi gửi các cấu xạ (tất cả đều đến cấu xạ $\operatorname{id}_1$).\hfill$\diamond$

Đôi khi ta gọi $\operatorname{Disc}(\underline{1})$ là \emph{phạm trù cuối} (xem Mục 6.1.3), và để đơn giản, ký hiệu nó là $\underline{1}$. Đối tượng duy nhất của nó được ký hiệu là $1$.
\paragraph{Bài tập 5.1.2.41.} Nếu ai đó nói: ``$\operatorname{Ob}$ là một hàm tử từ $\mathbf{Cat}$ đến $\mathbf{Set}$'', họ có thể muốn nói gì?\hfill$\diamond$
\paragraph{Lời giải 5.1.2.41.} Có lẽ họ muốn nói rằng có một hàm tử $\mathbf{Cat}\to\mathbf{Set}$ gửi một phạm trù $\mathcal C$ đến tập đối tượng $\operatorname{Ob}(\mathcal C)$ của nó. Vì người nói không cho biết hàm tử $\operatorname{Ob}$ này làm gì trên các cấu xạ, anh ấy ngụ ý rằng điều đó hiển nhiên. Một cấu xạ trong $\mathbf{Cat}$ là một hàm tử $F:\mathcal C\to\mathcal D$, mà theo định nghĩa, bao gồm một phần trên đối tượng. Nói cách khác, ý nghĩa của $\operatorname{Ob}(F):\operatorname{Ob}(\mathcal C)\to\operatorname{Ob}(\mathcal D)$ quả thật hiển nhiên, vì nó đã được cho trong dữ liệu xác định $F$ (xem Định nghĩa 5.1.2.1). Không khó kiểm tra rằng $\operatorname{Ob}$ bảo toàn đồng nhất và hợp thành, nên nó quả thật là một hàm tử.\hfill$\diamond$
% Source: printed p. 239; PDF p. 248
\paragraph{Bài tập 5.1.2.42.} Nếu ai đó nói: ``$\operatorname{Hom}$ là một hàm tử từ $\mathbf{Cat}$ đến $\mathbf{Set}$; ở đây $\operatorname{Hom}$ là ánh xạ gửi $\mathcal C$ đến tập $\operatorname{Hom}_{\mathcal C}$, như trong Bài tập 5.1.1.27'', họ có thể muốn nói gì?\hfill$\diamond$
\paragraph{Lời giải 5.1.2.42.} Có lẽ họ muốn nói rằng có một hàm tử $\mathbf{Cat}\to\mathbf{Set}$ gửi một phạm trù $\mathcal C$ đến tập cấu xạ $\operatorname{Hom}_{\mathcal C}$ của nó. Vì người nói không cho biết hàm tử $\operatorname{Hom}$ này làm gì trên các cấu xạ, cô ấy ngụ ý rằng điều đó hiển nhiên. Một cấu xạ trong $\mathbf{Cat}$ là một hàm tử $F:\mathcal C\to\mathcal D$, mà theo định nghĩa, bao gồm một phần trên cấu xạ. Nói cách khác, ý nghĩa của $\operatorname{Hom}(F):\operatorname{Hom}_{\mathcal C}\to\operatorname{Hom}_{\mathcal D}$ quả thật hiển nhiên, vì nó đã được cho trong dữ liệu xác định $F$ (xem Định nghĩa 5.1.2.1). Dễ kiểm tra rằng $\operatorname{Hom}$ bảo toàn đồng nhất và hợp thành, nên nó quả thật là một hàm tử.\hfill$\diamond$
\section{Các phạm trù và hàm tử thường gặp\texorpdfstring{\\}{ }trong toán học thuần túy}
\subsection{Monoid, nhóm, tiền thứ tự và đồ thị}
Ta đã thấy trong Mục 5.1.1 rằng có một phạm trù monoid $\mathbf{Mon}$, một phạm trù nhóm $\mathbf{Grp}$, một phạm trù tiền thứ tự $\mathbf{PrO}$, và một phạm trù đồ thị $\mathbf{Grph}$. Mục này cho thấy mỗi monoid $\mathcal M$, mỗi nhóm $\mathcal G$, và mỗi tiền thứ tự $\mathcal P$ đều có thể được xem là một phạm trù riêng. Nếu mỗi đối tượng trong $\mathbf{Mon}$ là một phạm trù, ta có thể hy vọng rằng mỗi cấu xạ trong $\mathbf{Mon}$ chỉ là một hàm tử; điều đó quả thật đúng. Điều tương tự cũng đúng với $\mathbf{Grp}$ và $\mathbf{PrO}$. Ta đã thấy trong Ví dụ 5.1.2.33 cách xem mỗi đồ thị như cho ra một phạm trù tự do. Một góc nhìn khác về đồ thị (tức là đồ thị như hàm tử) được bàn trong Mục 5.2.1.21.
\paragraph{5.2.1.1 Monoid như phạm trù.} Ví dụ 4.1.2.9 nói rằng để biểu diễn một monoid bằng olog, ta chỉ nên dùng một hộp. Ví dụ 4.5.3.3 lại nói rằng một tác động monoid có thể được nắm bắt bằng chỉ một bảng. Các ý tưởng này được gói gọn trong nhận thức rằng một monoid được mô hình hóa hoàn hảo bởi một phạm trù có một đối tượng.
\paragraph{Mỗi monoid như một phạm trù có một đối tượng.} Cho $(M,e,\star)$ là một monoid. Ta xem nó là một phạm trù $\mathcal M$ có một đối tượng, $\operatorname{Ob}(\mathcal M)=\{\blacktriangle\}$, và
\[\operatorname{Hom}_{\mathcal M}(\blacktriangle,\blacktriangle):=M.\]
% Source: printed p. 240; PDF p. 249
Cấu xạ đồng nhất $\operatorname{id}_{\blacktriangle}$ đóng vai trò đồng nhất $e$ của monoid, và công thức hợp thành
\[\circ:\operatorname{Hom}_{\mathcal M}(\blacktriangle,\blacktriangle)\times\operatorname{Hom}_{\mathcal M}(\blacktriangle,\blacktriangle)\longrightarrow\operatorname{Hom}_{\mathcal M}(\blacktriangle,\blacktriangle)\]
được cho bởi $\star:M\times M\to M$. Các luật kết hợp và đồng nhất của monoid khớp chính xác với các luật kết hợp và đồng nhất của phạm trù.

Nếu một monoid là một phạm trù có một đối tượng, liệu có cách nào diễn đạt khái niệm đồng cấu monoid bằng ngôn ngữ phạm trù không? Giả sử $\mathcal M=(M,e,\star)$ và $\mathcal M'=(M',e',\star')$. Ta biết rằng một đồng cấu monoid là một hàm $f:M\to M'$ sao cho $f(e)=e'$ và với mỗi cặp $m_0,m_1\in M$, ta có $f(m_0\star m_1)=f(m_0)\star' f(m_1)$. Một hàm tử $\mathcal M\to\mathcal M'$ là gì?
\paragraph{Mỗi đồng cấu monoid như một hàm tử giữa các phạm trù có một đối tượng.} Giả sử $\operatorname{Ob}(\mathcal M)=\{\blacktriangle\}$ và $\operatorname{Ob}(\mathcal M')=\{\blacktriangle'\}$; ta biết rằng $\operatorname{Hom}_{\mathcal M}(\blacktriangle,\blacktriangle)=M$ và $\operatorname{Hom}_{\mathcal M'}(\blacktriangle',\blacktriangle')=M'$. Một hàm tử $F:\mathcal M\to\mathcal M'$ trước hết gồm một hàm $\operatorname{Ob}(\mathcal M)\to\operatorname{Ob}(\mathcal M')$, nhưng mỗi tập này chỉ có một phần tử, nên ở phương diện đó không có gì để nói: ta phải có $F(\blacktriangle)=\blacktriangle'$. Nó còn gồm một hàm $\operatorname{Hom}_{\mathcal M}\to\operatorname{Hom}_{\mathcal M'}$, nhưng đó chỉ là một hàm $M\to M'$. Các công thức đồng nhất và hợp thành của hàm tử khớp chính xác với các công thức đồng nhất và hợp thành của đồng cấu monoid. Vì thế, một đồng cấu monoid chỉ là một hàm tử giữa các phạm trù có một đối tượng.
\paragraph{Khẩu hiệu 5.2.1.2.}\begin{quote}Một monoid là một phạm trù có một đối tượng. Một đồng cấu monoid chỉ là một hàm tử giữa các phạm trù có một đối tượng.\end{quote}
Điều này được phát biểu chính xác trong định lý sau.
\paragraph{Định lý 5.2.1.3.} Có một hàm tử $i:\mathbf{Mon}\to\mathbf{Cat}$ với các tính chất sau:
\begin{itemize}
\item Với mỗi monoid $\mathcal M\in\operatorname{Ob}(\mathbf{Mon})$, bản thân phạm trù $i(\mathcal M)\in\operatorname{Ob}(\mathbf{Cat})$ có đúng một đối tượng,
\[|\operatorname{Ob}(i(\mathcal M))|=1.\]
\item Với mỗi cặp monoid $\mathcal M,\mathcal M'\in\operatorname{Ob}(\mathbf{Mon})$, hàm
\[\operatorname{Hom}_{\mathbf{Mon}}(\mathcal M,\mathcal M')\xrightarrow{\cong}\operatorname{Hom}_{\mathbf{Cat}}(i(\mathcal M),i(\mathcal M')),\]
do hàm tử $i$ cảm sinh, là một song ánh.
\end{itemize}
% Source: printed p. 241; PDF p. 250
\begin{proof}Về cơ bản, đây chính là nội dung các đoạn trước. Hàm tử $i$ gửi một monoid đến phạm trù có một đối tượng tương ứng, và gửi một đồng cấu monoid đến hàm tử tương ứng. Có thể kiểm tra rằng $i$ bảo toàn đồng nhất và hợp thành.\end{proof}
Định lý 5.2.1.3 đặt lý thuyết monoid vào thế giới phạm trù một cách rất đẹp. Nhưng ta còn có những cách khác để nghĩ về monoid, cụ thể là tác động của chúng trên các tập hợp. Nếu có thể đưa cả tác động monoid vào lý thuyết phạm trù, câu chuyện sẽ vững vàng hơn nhiều; và ta có thể làm được điều đó.
\paragraph{Mỗi tác động monoid như một hàm tử giá trị tập hợp.} Nhắc lại từ Định nghĩa 4.1.2.1 rằng nếu $(M,e,\star)$ là một monoid, một tác động gồm một tập $S$ và một hàm $\circlearrowright:M\times S\to S$ sao cho $e\circlearrowright s=s$ và $m_0\circlearrowright(m_1\circlearrowright s)=(m_0\star m_1)\circlearrowright s$ với mọi $s\in S$. Ta có thể liên hệ khái niệm tác động monoid với khái niệm hàm tử như thế nào? Vì monoid tác động trên tập hợp, một ý tưởng là thử hỏi một hàm tử $F:\mathcal M\to\mathbf{Set}$ là gì; ý tưởng này sẽ có hiệu quả.

Monoid được xem như phạm trù $\mathcal M$ chỉ có một đối tượng $\blacktriangle$, nên $F$ cung cấp một tập $S:=F(\blacktriangle)\in\operatorname{Ob}(\mathbf{Set})$. Nó còn cung cấp một hàm $\operatorname{Hom}_F:\operatorname{Hom}_{\mathcal M}(\blacktriangle,\blacktriangle)\to\operatorname{Hom}_{\mathbf{Set}}(F(\blacktriangle),F(\blacktriangle))$, hay ngắn gọn hơn, một hàm
\[H_F:M\longrightarrow\operatorname{Hom}_{\mathbf{Set}}(S,S).\]
Bằng phép curry hóa (xem Mệnh đề 3.4.2.3), đây cũng chính là một hàm $\circlearrowright:M\times S\to S$. Luật tác động monoid thứ nhất, $e\circlearrowright s=s$, trở thành luật hàm tử bảo toàn đồng nhất, $\operatorname{Hom}_F(\operatorname{id}_{\blacktriangle})=\operatorname{id}_S$. Luật tác động monoid còn lại tương đương với luật hợp thành của hàm tử.
\paragraph{5.2.1.4 Nhóm như phạm trù.} Một nhóm chỉ là một monoid $(M,e,\star)$ mà mọi phần tử $m\in M$ đều khả nghịch, nghĩa là tồn tại $m'\in M$ sao cho $m\star m'=e=m'\star m$. Nếu một monoid cũng chính là một phạm trù $\mathcal M$ có một đối tượng, thì một nhóm phải là một phạm trù có một đối tượng, với một tính chất bổ sung liên quan đến tính khả nghịch. Các phần tử của $M$ là các cấu xạ của phạm trù $\mathcal M$, nên ta cần một khái niệm khả nghịch cho cấu xạ. May thay, ta đã có khái niệm đó: đẳng cấu.
\paragraph{Khẩu hiệu 5.2.1.5.}\begin{quote}Một nhóm là một phạm trù $\mathcal G$ có một đối tượng, sao cho mọi cấu xạ trong $\mathcal G$ đều là đẳng cấu. Một đồng cấu nhóm chỉ là một hàm tử giữa các phạm trù như vậy.\end{quote}
% Source: printed p. 242; PDF p. 251
\paragraph{Định lý 5.2.1.6.} Có một hàm tử $i:\mathbf{Grp}\to\mathbf{Cat}$ với các tính chất sau:
\begin{itemize}
\item Với mỗi nhóm $\mathcal G\in\operatorname{Ob}(\mathbf{Grp})$, bản thân phạm trù $i(\mathcal G)\in\operatorname{Ob}(\mathbf{Cat})$ có đúng một đối tượng, và mọi cấu xạ $m$ trong $i(\mathcal G)$ đều là đẳng cấu.
\item Với mỗi cặp nhóm $\mathcal G,\mathcal G'\in\operatorname{Ob}(\mathbf{Grp})$, hàm
\[\operatorname{Hom}_{\mathbf{Grp}}(\mathcal G,\mathcal G')\xrightarrow{\cong}\operatorname{Hom}_{\mathbf{Cat}}(i(\mathcal G),i(\mathcal G')),\]
do hàm tử $i$ cảm sinh, là một song ánh.
\end{itemize}
Cũng như với monoid, một tác động của một nhóm $(G,e,\star)$ trên một tập $S\in\operatorname{Ob}(\mathbf{Set})$ cũng chính là một hàm tử $\mathcal G\to\mathbf{Set}$ gửi đối tượng duy nhất của $\mathcal G$ đến tập $S$.
\paragraph{5.2.1.7 Một monoid và một nhóm đặt tại mỗi đối tượng trong một phạm trù bất kỳ.} Nếu một monoid chỉ là một phạm trù có một đối tượng, ta có thể tìm các monoid trong bất kỳ phạm trù $\mathcal C$ nào bằng cách tập trung vào một đối tượng của $\mathcal C$. Điều tương tự cũng đúng với nhóm.
\paragraph{Ví dụ 5.2.1.8 (Monoid tự cấu xạ).} Cho $\mathcal C$ là một phạm trù và $x\in\operatorname{Ob}(\mathcal C)$ là một đối tượng. Đặt $M=\operatorname{Hom}_{\mathcal C}(x,x)$. Chú ý rằng với hai phần tử bất kỳ $f,g\in M$, ta có $f\circ g:x\to x$ trong $M$. Đặt $\mathcal M=(M,\operatorname{id}_x,\circ)$. Dễ kiểm tra rằng $\mathcal M$ là một monoid; nó được gọi là \emph{monoid tự cấu xạ} của $x$ trong $\mathcal C$, ký hiệu $\operatorname{End}(x)$.
\paragraph{Ví dụ 5.2.1.9 (Nhóm tự đẳng cấu).} Cho $\mathcal C$ là một phạm trù và $x\in\operatorname{Ob}(\mathcal C)$ là một đối tượng. Đặt $G=\{f\in\operatorname{Hom}_{\mathcal C}(x,x)\mid f\text{ là một đẳng cấu}\}$. Đặt $\mathcal G=(G,\operatorname{id}_x,\circ)$. Có thể kiểm tra rằng $\mathcal G$ là một nhóm; nó được gọi là \emph{nhóm tự đẳng cấu} của $x$ trong $\mathcal C$, ký hiệu $\operatorname{Aut}(x)$.
\paragraph{Bài tập 5.2.1.10.} Cho $S=\{1,2,3,4\}\in\operatorname{Ob}(\mathbf{Set})$.
\begin{enumerate}[label=\alph*.]
\item Nhóm tự đẳng cấu $\operatorname{Aut}(S)$ của $S$ trong $\mathbf{Set}$ là gì, và nhóm này có bao nhiêu phần tử?
\item Monoid tự cấu xạ $\operatorname{End}(S)$ của $S$ trong $\mathbf{Set}$ là gì, và monoid này có bao nhiêu phần tử?
\item Nhắc lại từ Ví dụ 5.1.2.3 rằng mỗi nhóm đều có một monoid nền $U(G)$. Monoid tự cấu xạ của $S$ có phải là monoid nền của nhóm tự đẳng cấu của $S$ không? Tức là có phải $\operatorname{End}(S)=U(\operatorname{Aut}(S))$ không?
\end{enumerate}\hfill$\diamond$
% Source: printed p. 243; PDF p. 252
\paragraph{Lời giải 5.2.1.10.}
\begin{enumerate}[label=\alph*.]
\item Đó là nhóm hoán vị $\operatorname{Aut}S=(\operatorname{Iso}(S),\operatorname{id}_S,\circ)$, trong đó $\operatorname{Iso}(S)$ được định nghĩa trong Bài tập 4.2.1.7. Nó có $4!=24$ phần tử.
\item Đó là monoid $(\operatorname{Hom}_{\mathbf{Set}}(S,S),\operatorname{id}_S,\circ)$ của mọi hàm từ $S$ đến $S$. Nó có $4^4=256$ phần tử.
\item Không. Monoid nền của một nhóm $G$ có cùng số phần tử với $G$, nên monoid nền của nhóm hoán vị của $S$ không phải là monoid tự cấu xạ của $S$.

Bạn đọc có thể nhận thấy có một hàm tử $\operatorname{Core}:\mathbf{Mon}\to\mathbf{Grp}$ nhận một monoid $M$ và trả về tập con các phần tử khả nghịch của nó; tập con này tạo thành một nhóm gọi là \emph{lõi} của $M$. Lõi của monoid tự cấu xạ $\operatorname{End}(S)$ quả thật là nhóm tự đẳng cấu $\operatorname{Core}(\operatorname{End}(S))=\operatorname{Aut}(S)$.
\end{enumerate}\hfill$\diamond$
\paragraph{Bài tập 5.2.1.11.} Xét đồ thị $G$ sau, có bốn đỉnh và tám mũi tên:
\begin{center}\begin{tikzpicture}[ctfivearrow]
\node[circle,fill,inner sep=1.4pt,label=above:$1$] (a) at (0,2) {};
\node[circle,fill,inner sep=1.4pt,label=above:$2$] (b) at (3,2) {};
\node[circle,fill,inner sep=1.4pt,label=below:$3$] (c) at (0,0) {};
\node[circle,fill,inner sep=1.4pt,label=below:$4$] (d) at (3,0) {};
\draw (a) to[bend left=10] node[above]{$12$} (b);
\draw (b) to[bend left=10] node[below]{$21$} (a);
\draw (a) to[bend left=14] node[right]{$13$} (c);
\draw (c) to[bend left=14] node[left]{$31$} (a);
\draw (b) to[bend left=14] node[right]{$24$} (d);
\draw (d) to[bend left=14] node[left]{$42$} (b);
\draw (c) to[bend left=10] node[above]{$34$} (d);
\draw (d) to[bend left=10] node[below]{$43$} (c);
\end{tikzpicture}\end{center}
Nhóm tự đẳng cấu $\operatorname{Aut}(G)$ của $G\in\operatorname{Ob}(\mathbf{Grph})$ là gì? Gợi ý: Mỗi tự đẳng cấu của $G$ sẽ cảm sinh một tự đẳng cấu của tập $\{1,2,3,4\}$; những tự đẳng cấu nào bảo toàn các đầu mút của mũi tên?\hfill$\diamond$
\paragraph{Lời giải 5.2.1.11.} Ta dùng nhận thức thị giác để định hướng. Đồ thị $G$ có hình dạng một hình vuông. Trong $4!$ tự đẳng cấu khác nhau có thể có của $\{1,2,3,4\}$, chỉ những tự đẳng cấu bảo toàn hình dạng hình vuông mới là tự đẳng cấu của $G$. Nhóm tự đẳng cấu của $G$ được gọi là \emph{nhóm nhị diện cấp 8} (xem Ví dụ 4.2.1.4). Nó có tám phần tử,
\[\{e,r,r^2,r^3,f,fr,fr^2,fr^3\},\]
% Source: printed p. 244; PDF p. 253
trong đó $r$ nghĩa là xoay hình vuông $90^\circ$ theo chiều kim đồng hồ, và $f$ nghĩa là lật hình vuông theo chiều ngang. Chẳng hạn, có thể lật hình vuông theo chiều dọc bằng cách lật theo chiều ngang rồi xoay hai lần: $fr^2$.\hfill$\diamond$
\paragraph{5.2.1.12 Tiền thứ tự như phạm trù.} Một tiền thứ tự $(X,\leq)$ gồm một tập $X$ và một quan hệ hai ngôi $\leq$ có tính phản xạ và bắc cầu. Từ $(X,\leq)\in\operatorname{Ob}(\mathbf{PrO})$, ta có thể tạo ra một phạm trù $\mathcal X\in\operatorname{Ob}(\mathbf{Cat})$ như sau. Định nghĩa $\operatorname{Ob}(\mathcal X)=X$, và với hai đối tượng bất kỳ $x,y\in X$, định nghĩa
\[\operatorname{Hom}_{\mathcal X}(x,y)=\begin{cases}\{\text{``}x\leq y\text{''}\}&\text{nếu }x\leq y,\\\varnothing&\text{nếu }x\nleq y.\end{cases}\]
Nói rõ hơn: nếu $x\leq y$, ta đặt $\operatorname{Hom}_{\mathcal X}(x,y)$ là tập chứa đúng một phần tử, cụ thể là chuỗi ``$x\leq y$''.\footnote{Tên của cấu xạ này không quan trọng. Điều quan trọng là $\operatorname{Hom}_{\mathcal X}(x,y)$ có đúng một phần tử khi và chỉ khi $x\leq y$.} Nếu cặp $(x,y)$ không thuộc quan hệ $\leq$, ta đặt $\operatorname{Hom}_{\mathcal X}(x,y)$ là tập rỗng. Công thức hợp thành
\begin{equation*}\circ:\operatorname{Hom}_{\mathcal X}(x,y)\times\operatorname{Hom}_{\mathcal X}(y,z)\longrightarrow\operatorname{Hom}_{\mathcal X}(x,z)\tag{5.7}\end{equation*}
được xác định hoàn toàn, vì một trong hai khả năng xảy ra. Khả năng thứ nhất là vế trái rỗng (nếu $x\nleq y$ hoặc $y\nleq z$; trong trường hợp này có duy nhất một hàm $\circ$ như trong (5.7)). Khả năng còn lại là vế trái không rỗng, trong trường hợp $x\leq y$ và $y\leq z$, suy ra $x\leq z$; vì vậy, vế phải có đúng một phần tử ``$x\leq z$'', và một lần nữa có duy nhất một hàm $\circ$ như trong (5.7).

Mặt khác, nếu $\mathcal C$ là một phạm trù có tính chất rằng với mỗi cặp đối tượng $x,y\in\operatorname{Ob}(\mathcal C)$, tập $\operatorname{Hom}_{\mathcal C}(x,y)$ hoặc rỗng hoặc có một phần tử, thì ta có thể tạo một tiền thứ tự từ $\mathcal C$. Cụ thể, đặt $X=\operatorname{Ob}(\mathcal C)$ và nói $x\leq y$ nếu có một cấu xạ $x\to y$ trong $\mathcal C$.
\paragraph{Mệnh đề 5.2.1.13.} Có một hàm tử $i:\mathbf{PrO}\to\mathbf{Cat}$ với các tính chất sau cho mỗi tiền thứ tự $(X,\leq)$:
\begin{enumerate}[label=\arabic*.]
\item Phạm trù $\mathcal X:=i(X,\leq)$ có các đối tượng $\operatorname{Ob}(\mathcal X)=X$.
\item Với mỗi cặp phần tử $x,x'\in\operatorname{Ob}(\mathcal X)$, tập $\operatorname{Hom}_{\mathcal X}(x,x')$ có nhiều nhất một phần tử.
\end{enumerate}
Hơn nữa, mọi phạm trù có tính chất 2 đều thuộc ảnh của hàm tử $i$.
% Source: printed p. 245; PDF p. 254
\begin{proof}Để xác định một hàm tử $i:\mathbf{PrO}\to\mathbf{Cat}$, ta cần nói nó làm gì trên đối tượng và trên cấu xạ. Với một đối tượng $(X,\leq)$ trong $\mathbf{PrO}$, ta gán phạm trù $\mathcal X$ có các đối tượng $X$ và một cấu xạ duy nhất $x\to x'$ nếu $x\leq x'$. Với một cấu xạ tiền thứ tự $f:(X,\leq_X)\to(Y,\leq_Y)$, ta phải gán một hàm tử $i(f):\mathcal X\to\mathcal Y$. Một lần nữa, để xác định một hàm tử, ta cần nói nó làm gì trên các đối tượng và cấu xạ của $\mathcal X$. Với một đối tượng $x\in\operatorname{Ob}(\mathcal X)=X$, ta gán đối tượng $f(x)\in Y=\operatorname{Ob}(\mathcal Y)$. Cho một cấu xạ $f:x\to x'$ trong $\mathcal X$, ta biết rằng $x\leq x'$, nên theo Định nghĩa 4.4.4.1, ta có $f(x)\leq f(x')$; ta gán cho $f$ cấu xạ duy nhất $f(x)\to f(x')$ trong $\mathcal Y$. Việc kiểm tra rằng các quy tắc của hàm tử (bảo toàn đồng nhất và hợp thành) được tuân thủ là thông thường.\end{proof}
\paragraph{Khẩu hiệu 5.2.1.14.}\begin{quote}Một tiền thứ tự là một phạm trù mà mỗi tập hom có 0 hoặc 1 phần tử. Một cấu xạ tiền thứ tự chỉ là một hàm tử giữa các phạm trù như vậy.\end{quote}
\paragraph{Bài tập 5.2.1.15.} Giả sử $\mathcal C$ là một tiền thứ tự (được xem như phạm trù). Cho $x,y\in\operatorname{Ob}(\mathcal C)$ là các đối tượng sao cho $x\leq y$ và $y\leq x$. Chứng minh rằng có một đẳng cấu $x\to y$ trong $\mathcal C$.\hfill$\diamond$
\paragraph{Lời giải 5.2.1.15.} Theo ngôn ngữ phạm trù, $x\leq y$ nghĩa là có một cấu xạ $f:x\to y$ trong $\mathcal C$, và tương tự, $g:y\leq x$ nghĩa là có một cấu xạ $y\to x$ trong $\mathcal C$. Ta có thể hợp thành chúng để được một cấu xạ $g\circ f:x\to x$ và một cấu xạ $f\circ g:y\to y$. Nhưng điều kiện để $\mathcal C$ là một tiền thứ tự là mỗi tập hom có nhiều nhất một phần tử. Ta đã có $\operatorname{id}_x:x\to x$ và $\operatorname{id}_y:y\to y$, nên các hợp thành $g\circ f$ và $f\circ g$ phải lần lượt là các cấu xạ đồng nhất $\operatorname{id}_x$ và $\operatorname{id}_y$. Vì vậy, theo định nghĩa, $f$ và $g$ là đẳng cấu.\hfill$\diamond$
\paragraph{Bài tập 5.2.1.16.} Mệnh đề 5.2.1.13 phát biểu rằng một tiền thứ tự có thể được xem là một phạm trù $\mathcal P$. Nhắc lại từ Định nghĩa 4.4.1.1 rằng một thứ tự bộ phận là một tiền thứ tự có thêm một tính chất. Hãy diễn đạt tính chất định nghĩa thứ tự bộ phận bằng đẳng cấu trong phạm trù $\mathcal P$.\hfill$\diamond$
\paragraph{Lời giải 5.2.1.16.} Một tiền thứ tự là một phạm trù $\mathcal P$ sao cho với mỗi cặp đối tượng $x,y\in\operatorname{Ob}(\mathcal P)$, tập $\operatorname{Hom}_{\mathcal P}(x,y)$ có nhiều nhất một phần tử. Một phần tử của tập hom này biểu thị sự kiện
